\documentclass[11pt,aspectratio=169]{beamer}

\usepackage{amsmath,amssymb,bm}
\usepackage{booktabs}
\usepackage{tikz}
\usetikzlibrary{arrows.meta,positioning,calc}

\usetheme{default}
\setbeamertemplate{navigation symbols}{}
\setbeamertemplate{headline}{}
\setbeamertemplate{footline}{%
  \leavevmode
  \hbox{%
    \begin{beamercolorbox}[wd=\paperwidth,ht=2.6ex,dp=1.1ex,right]{page number in head/foot}%
      \insertframenumber/\inserttotalframenumber\hspace{0.35cm}
    \end{beamercolorbox}%
  }%
}
\setbeamertemplate{frametitle}{\vspace{0.25cm}\insertframetitle\par\vspace{0.08cm}\hrule}
\setbeamercolor{normal text}{fg=black,bg=white}
\setbeamercolor{frametitle}{fg=black,bg=white}
\setbeamercolor{structure}{fg=black}
\setbeamersize{text margin left=0.7cm,text margin right=0.7cm}

\newcommand{\dd}{\,\mathrm{d}}
\newcommand{\Om}{\Omega}
\newcommand{\GamD}{\Gamma_D}
\newcommand{\GamN}{\Gamma_N}
\newcommand{\Th}{\mathcal T_h}
\newcommand{\Fh}{\mathcal F_h}
\newcommand{\Vzero}{\mathcal V_0}
\newcommand{\Vhzero}{\mathcal V_{h,0}}
\newcommand{\Qh}{\mathcal Q_h}
\newcommand{\Ih}{\mathcal I_h}
\newcommand{\norm}[2][]{\left\lVert#2\right\rVert_{#1}}
\newcommand{\seminorm}[2][]{\left\lvert#2\right\rvert_{#1}}
\newcommand{\grad}{\nabla}
\newcommand{\divg}{\nabla\!\cdot}
\newcommand{\bn}{\bm n}
\newcommand{\bq}{\bm q}
\newcommand{\bx}{\bm x}

\begin{document}


\begin{frame}[plain]
  \centering\vfill
  {\LARGE\bfseries Interpretation and Adaptive Refinement\par}
  \vspace{0.45cm}
  {\large Part 3: independent study\par}
  \vspace{0.25cm}
  {\large Efficiency, effectivity, marking, and the adaptive loop\par}
  \vfill
\end{frame}

\begin{frame}{Reading outline}
Part 2 derived a computable estimator and proved
\[
 \norm[a]{u-u_h}\leq C_{\mathrm{rel}}\eta,
 \qquad \eta^2=\sum_e\eta_e^2.
\]
This part addresses how to interpret and use these numbers.
\begin{enumerate}
 \item Distinguish reliability, local efficiency, and effectivity.
 \item Select elements using their local indicators.
 \item Refine the selected mesh and compute a new FE solution.
 \item Interpret a stopping test and recognize what it does not guarantee.
\end{enumerate}
The discussion continues the scalar-diffusion setting of Parts 1 and 2.
\end{frame}

\begin{frame}{Reliability and local efficiency}
Reliability bounds the global energy error from above:
\[
 \norm[a]{e}\leq C_{\mathrm{rel}}\eta,
 \qquad e=u-u_h.
\]
Local efficiency asks for a converse comparison on a surrounding patch:
\[
 \eta_e\leq C_{\mathrm{eff}}
 \bigl(\norm[a,\omega_e]{e}+\mathrm{osc}_e\bigr),
 \qquad
 \norm[a,\omega_e]{e}^2=\int_{\omega_e}k\grad e\cdot\grad e\,\dd x.
\]
Here $\mathrm{osc}_e$ measures unresolved variation in the prescribed data.
The constants must be independent of mesh size on the stated mesh family.

Reliability makes a small $\eta$ meaningful. When data oscillation is small,
local efficiency ties a large indicator to error near that element.

\vspace{0.12cm}
{\footnotesize Chapter 8, \S8.5 defines local efficiency. Neither that chapter
nor this deck proves the displayed efficiency bound for the general problem.}
\end{frame}

\begin{frame}{Meaning of data oscillation}
Consider the simpler case $k=1$ and $u=0$ on the whole boundary.
An example of a data-oscillation term uses the element average
\[
 \overline f_e=\frac{1}{|\mathcal K_e|}\int_{\mathcal K_e}f\,\dd x,
 \qquad
 \mathrm{osc}_e=h_e\norm[L^2(\mathcal K_e)]{f-\overline f_e}.
\]
The difference $f-\overline f_e$ is the part of the source not represented
by that constant average. For constant $f$ on an element, this term is zero.

With Neumann data, a data-oscillation term can similarly compare $q$ with
its local polynomial approximation on a boundary facet.
The worked examples here use constant source and Neumann data on each
element or facet, so these oscillation terms vanish.
\end{frame}

\begin{frame}{The effectivity index}
For a verification problem with known exact solution and nonzero error, define
\[
 \boxed{I_{\mathrm{eff}}=\frac{\eta}{\norm[a]{u-u_h}}.}
\]
\begin{itemize}
 \item $I_{\mathrm{eff}}=1$ means that estimator and energy error agree.
 \item $I_{\mathrm{eff}}>1$ means that the estimator exceeds the error by that factor.
 \item Reliability gives $I_{\mathrm{eff}}\geq1/C_{\mathrm{rel}}$.
 With an unknown $C_{\mathrm{rel}}$, this does not imply $I_{\mathrm{eff}}\geq1$.
\end{itemize}
Exact error is unavailable in a general computation, so effectivity is a
verification measure, not a stopping test for an unknown solution.

\vspace{0.15cm}
{\footnotesize Chapter 8, equation for the effectivity index and
``Effectivity for the uniformly loaded bar.''}
\end{frame}

\begin{frame}{Effectivity in the nonuniform 1D example}
Part 2 obtained the squared estimator contributions and squared element errors
\[
 \left(\eta_{1,\mathrm{1D}}^2,\eta_{2,\mathrm{1D}}^2\right)
 =\left(\frac1{27},\frac8{27}\right),\qquad
 \left(E_1^2,E_2^2\right)=\left(\frac1{324},\frac8{324}\right).
\]
Here $E_e^2=\int_{\mathcal K_e}|u'-u_h'|^2\,\dd x$ and
$\eta_{e,\mathrm{1D}}^2=h_e^3$. Summing gives
\[
 \eta_{\mathrm{1D}}^2=\frac13,\qquad
 \norm[a]{e}^2=\frac1{36},\qquad
 I_{\mathrm{eff}}=\frac{1/\sqrt3}{1/6}=\sqrt{12}\approx3.46.
\]
The explicit 1D bound $\norm[a]{e}\leq\eta_{\mathrm{1D}}$ is valid, but
the estimator is about $3.46$ times the actual error in this example.

This calculation checks one example. It does not establish a general
efficiency theorem.
\end{frame}

\begin{frame}{Marking elements for refinement}
Let $\mathcal M\subset\Th$ be the set of elements selected for refinement.
The indicators determine this set before a new mesh is constructed.

\textbf{Maximum marking.} Choose $0<\gamma<1$ and mark every element with
\[
 \eta_e\geq\gamma\max_j\eta_j.
\]
\textbf{D\"orfler (bulk) marking.} Choose $0<\theta<1$ and select enough
elements that
\[
 \boxed{\sum_{\mathcal K_e\in\mathcal M}\eta_e^2
 \geq\theta\sum_e\eta_e^2=\theta\eta^2.}
\]
A simple procedure sorts the squared indicators from largest to smallest
and includes them until the required fraction is reached.

\vspace{0.12cm}
{\footnotesize Arnold, Chapter 4, \S9.4, printed p.~73 (PDF p.~77).
Here $\theta$ denotes the fraction of the squared total.}
\end{frame}

\begin{frame}{Bulk marking in the 1D example}
Use the 1D estimator from Part 2, not the multidimensional facet formula:
\[
 \eta_{1,\mathrm{1D}}^2=\frac1{27},\qquad
 \eta_{2,\mathrm{1D}}^2=\frac8{27},\qquad
 \eta_{\mathrm{1D}}^2=\frac13.
\]
Take $\theta=0.8$. The second element contributes the fraction
\[
 \frac{\eta_{2,\mathrm{1D}}^2}{\eta_{\mathrm{1D}}^2}
 =\frac{8/27}{9/27}=\frac89>0.8.
\]
Therefore $\mathcal M=\{\mathcal K_2\}$ meets the bulk criterion.

The source residual is $1$ on both elements. The second element is marked
because its larger length produces the larger weighted contribution.
\end{frame}

\begin{frame}{Refining the marked element}
Bisect $\mathcal K_2=(1/3,1)$ at $x=2/3$. The new mesh is
\[
 0,\quad 1/3,\quad 2/3,\quad 1,
\]
so every new element has length $1/3$.
Solve the FE problem again on this mesh. For $P_1$ elements, $k=1$, and
$f=1$, the new element residuals are still $r_e=1$. Hence
\[
 \eta_{\mathrm{1D,new}}^2
 =3\left(\frac13\right)^3=\frac19
 <\frac13=\eta_{\mathrm{1D,old}}^2.
\]
This reduction was calculated for this problem. It does not guarantee the
same reduction for every problem or marked mesh.

In 2D and 3D, refinement must also preserve conformity and shape regularity.
Neighboring elements may need refinement to avoid hanging nodes.
\end{frame}

\begin{frame}{The adaptive finite element loop}
Begin with a conforming, shape-regular mesh.
\begin{enumerate}
 \item \textbf{Solve:} compute $u_h$ on the current mesh.
 \item \textbf{Estimate:} evaluate the element and facet residuals, then
 $\eta_e$ and $\eta=(\sum_e\eta_e^2)^{1/2}$.
 \item \textbf{Check:} stop if the chosen estimator tolerance is met or
 a stated resource limit is reached.
 \item \textbf{Mark:} select elements using the local indicators.
 \item \textbf{Refine:} construct a conforming new mesh and return to Step 1.
\end{enumerate}
Refining the mesh changes the FE space and the computed solution. Recompute
the indicators after each new solve rather than retaining the old values.

\vspace{0.15cm}
{\footnotesize Arnold, \S9.4, p.~73. Larson--Bengzon, \S4.9.2.}
\end{frame}

\begin{frame}{Stopping criteria and the unknown constant}
Given an estimator tolerance $\varepsilon_{\mathrm{est}}>0$, a common test is
\[
 \eta\leq\varepsilon_{\mathrm{est}}.
\]
Reliability then gives
\[
 \norm[a]{u-u_h}\leq C_{\mathrm{rel}}\varepsilon_{\mathrm{est}}.
\]
If $C_{\mathrm{rel}}$ is unknown, this test does not certify the same numerical
tolerance for the true error.

For the 1D example, the proved bound has $C_{\mathrm{rel}}=1$. There,
$\eta_{\mathrm{1D}}\leq\varepsilon$ does guarantee
$\norm[a]{u-u_h}\leq\varepsilon$.

A solver or mesh-size limit is a separate stopping reason. Reaching such a
limit does not establish that an error tolerance was achieved.
\end{frame}

\begin{frame}{Reading and practice}
\begin{enumerate}
 \item Read Part 2 through the estimator definition and reliability proof.
 Track where Galerkin orthogonality, each local estimate, and
 Cauchy--Schwarz enter the argument.
 \item Work through the 1D and two-triangle calculations in Part 2.
 Distinguish squared indicator contributions from the exact element errors.
 \item Read Part 3 through the marking calculation and adaptive loop.
 Apply the bulk criterion before deciding which elements to bisect.
\end{enumerate}
Book reading:
\href{https://ceadpx.github.io/finite-element-methods/chapters/08-aposteriori-adaptivity.html\#sec-c8-estimator}{Chapter 8, \S8.4},
followed by
\href{https://ceadpx.github.io/finite-element-methods/chapters/08-aposteriori-adaptivity.html\#sec-c8-reliability-efficiency}{\S8.5}.
The marking and adaptive-loop discussion is supplied by this deck and Arnold
\S9.4 because those sections of the book are not yet written.
\end{frame}

\begin{frame}{Scope of the residual-estimation module}
The three parts develop the following argument:
\begin{itemize}
 \item The variational residual equals the energy action of the unknown error.
 \item Element and facet residuals produce computable weighted indicators.
 \item Quasi-interpolation and finite patch overlap give a reliability bound.
 \item The indicators guide marking, and a new solve supplies the next estimate.
\end{itemize}
Efficiency and effectivity describe additional properties of an estimator.
A single computed effectivity index does not prove efficiency.

The general local-efficiency proof, convergence/optimality of the adaptive
algorithm, FEniCSx implementation, elasticity estimators, and goal-oriented
estimation are not developed in these decks.
\end{frame}

\begin{frame}{References and further reading}
\begin{itemize}
 \item Course book, Chapter 8, \S\S8.4--8.5: estimator, reliability,
 1D effectivity, and definitions of efficiency.
 \item D.~N.~Arnold, \emph{Lecture Notes on Numerical Analysis of Partial
 Differential Equations}, Chapter 4, \S9.4, printed p.~73 (PDF p.~77):
 local indicators, maximum and D\"orfler marking, and the adaptive loop.
 \item M.~G.~Larson and F.~Bengzon, \emph{The Finite Element Method:
 Theory, Implementation, and Practice}, \S2.6.1, pp.~38--39, and
 \S4.9, pp.~88--92: residual bounds and adaptive refinement.
\end{itemize}
For the operator and reliability proof used in Part 2:
\href{https://jschoeberl.github.io/iFEM/aposteriori/residualEE.html}{J.~Sch\"oberl, \emph{Interactive Finite Elements}, \S\S27.1--27.2.}
\end{frame}

\end{document}
